% 	
%	
%
\documentclass[a4paper, 12pt, twoside]{scrreprt}

% Encoding, islands characters.
%\usepackage{ucs}
\usepackage[utf8]{inputenc}
\usepackage[portuguese]{babel}
\usepackage{t1enc}
%\usepackage{csquotes} % Problemas com aspas na biblioteca.

\usepackage[intoc]{nomencl}
\usepackage{enumerate,color}
\usepackage{url}
\usepackage[pdfborder={0 0 0}]{hyperref}
\usepackage{appendix}
\usepackage{eso-pic}
\usepackage{amsmath}
\usepackage{amssymb}
\usepackage[nottoc]{tocbibind}



\usepackage[format=plain,labelformat=simple,labelsep=colon]{caption}
\usepackage{placeins}

% comprimidos
\usepackage{tabularx}

% Desenhos e fotos.
\usepackage{graphicx}
\usepackage[sf,normalsize]{subfigure}
\usepackage{tikz}
\usetikzlibrary{calc}
\usepackage{circuitikz}
\usepackage{tikz,pgffor}
\usetikzlibrary{calc,through,intersections, matrix}
\usepackage{tkz-euclide}
\usetkzobj{all}
\usepackage{geometry}

% units package
\usepackage{siunitx}
\usepackage{cancel}

% Suplemento Bibliográfico.
\usepackage[style=ieee, backend=biber]{biblatex} 
\addbibresource{heimildaskra.bib} % Fontes são baixadas a partir desta lista.

% Configurations
%\graphicspath{{figs/}}

%\usepackage{fancyhdr}
 
%\pagestyle{fancy}
%\fancyhf{}
%\rhead{Share\LaTeX}
%\lhead{Guides and tutorials}
%\rfoot{Page \thepage}


\setlength{\parskip}{\baselineskip}
\setlength{\parindent}{0cm}
\raggedbottom
% \setkomafont{subsection}{\normalfont\sffamily}

% Como modificação de fios para ser
% \setkomafont{captionlabel}{\itshape}
% \setkomafont{caption}{\itshape}

% Mais bela solução << esta é uma escolha vergonhosa
 \setkomafont{captionlabel}{\itshape}
 \setkomafont{caption}{\itshape}
 \setkomafont{section}{\FloatBarrier\Large}
 \setcapwidth[l]{\textwidth}
%\setcapindent{1em}


% Times new roman
%\usepackage[T1]{fontenc}
%\usepackage{mathptmx}

%
%	Preencha Informações sobre o titulo, departamento e autor.
%
\def\thesisyear{2016}       				% Ano Projecto Entrega
\def\thesismonth{Novembro}					% Mês do Projecto final
\def\theisnrofauthors{3}
\def\thesisauthors{Marina Camilo - up201307722 - up201307722@fe.up.pt
				\\ Diogo Ferreira - up201502853 - diogoff@fe.up.pt
                \\ Ângela Cardoso - up200204375 - angela.cardoso@fe.up.pt}						% Authors
\def\thesistitle{T05: Exploração de Marte} 		% Titulo
\def\thesisshorttitle{Sistemas Distribuidos} 	% Tiulo Curto

\def\thesisnroftutors{2}						% 1 = singular, >1 plural
\def\thesistutors{Eugénio da Costa Oliveira \\
				  Henrique Daniel de Avelar Lopes Cardoso }	% Formador

\begin{document}

\begin{titlepage}
\thispagestyle{empty}

\begin{tikzpicture}[remember picture,overlay]
  \coordinate [above=0cm] (toppoint) at (current page.north);
  % Logo Universidade.
  \node[anchor=north] (logonator) at ($(toppoint)+(0cm,-4.2cm)$) {


\includegraphics[width=8cm]{./FEUP_Logo.png}
    };

  % Titulo
  \node[anchor=north, align=center] (titlenator) at ($(logonator.south)+(0cm,-2.0cm)$) {
    \huge \sffamily \bfseries \thesistitle
    };

  \coordinate [above=0cm] (bottompoint) at (current page.south);

  % Promoçao de texto localização e atribuições
  \node[anchor=base, 
  		minimum width=\paperwidth, 
        align=center] at ($(bottompoint)+(0cm,5cm)$)  
    {
    \parbox{
      1\paperwidth
      }{
      \centering
      \bfseries\sffamily \Large 
	  4ºAno - Sistemas Distribuidos \\
	  \thesismonth~\thesisyear
      }
    };

  % nome do autor
  \node[anchor=north, align=center] (authornator) at ($(titlenator.south)+(0cm,-1.0cm)$) {
  	\parbox{
      1\paperwidth
      }{
      \centering
      \normalfont \Large \sffamily \thesisauthors
      }
    };
    
    %\normalfont \Large \sffamily \thesisauthors
	%};
    
  % nomes dos Supervisores
  \node[anchor=north, align=center] (advisornator) at ($(authornator.south)+(0cm,-1.0cm)$) {
	\parbox{
      1\paperwidth
      }{
      \centering
      \sffamily \Large Professores\\
      \thesistutors
      }
    };
    
\end{tikzpicture}

\end{titlepage}
\pagenumbering{arabic}

\setcounter{page}{2}
\newpage

\tableofcontents
%\listoffigures
%\listoftables

\clearpage
\chapter{Enunciado}
\pagenumbering{arabic}
\setcounter{page}{3}
%escrever aqui
\section{Descrição do cenário}
%escrever aqui
\section{Objectivos do trabalho}
%escrever aqui 
\section{Resultados esperados e forma de avaliação}
%escrever aqui
 
\chapter{Plataforma/Ferramenta}
%escrever aqui
\section{Para que serve}
%escrever aqui
\section{Descrição das características principais}
%escrever aqui
\section{Realce das funcionalidades relevantes para o trabalho}
%escrever aqui

\chapter{Especificação}
%escrever aqui
\section{Identificação e caracterização dos agentes (arquitectura, comportamento, estratégias)}
%escrever aqui
\section{Protocolos de interacção}
%escrever aqui
\section{Faseamento do projecto}
%escrever aqui

\chapter{Recursos}
%escrever aqui
\section{Bibliografia}
%escrever aqui
\section{Software}
%escrever aqui


\chapter{Exemplo Figura e tabela}

\begin{figure}[!htb]
\centering
\includegraphics[width=0.60\textwidth]{FEUP_Logo}
\caption[alguma coisa.]{Mais aluguma coisa} \label{fig:Array}
\end{figure}

\begin{table}[htb]
\centering
\caption{Alguma coisa.}
     \sffamily \begin{tabularx}{1.0\textwidth}{ p{5cm}  p{5cm}  p{5cm} }
    \hline
   \textbf{nome1} \hfill & \textbf{nome2} \hfill & \textbf{nome3} \\ \hline
    1.1 & 1.2 & 1.3\\
    2.1 & 2.2 & 2.3\\ \hline
    \end{tabularx} \normalfont
\label{table:Emissivity}
\end{table}

\section{Baixando a partir de outro documento}
Exemplos de autorizado~\cite{dirac} e tambem~\cite{einstein}.
%
% Skjal til að þætta inní aðalskjal
%
\section{Efni fengið frá ytra skjali}
Þetta efni er sótt úr skránni \emph{YtraEfniFyrirInclude.tex}. Hægt er að nota nokkrar skipanir til að þætta inn önnur skjöl.
\begin{itemize}
	\item {\textbackslash}input, færir inn texta án frekari vinnslu.
    \item {\textbackslash}include, færir inn texta 
		og setur {\textbackslash}clearpage í byrjun og enda texta.
    \item {\textbackslash}includeonly\{filename,filename2,...\}, skal setja
    	í {\glqq}Preamble{\grqq} til að velja hvaða skrár eru sótta með 
        {\textbackslash}include.
\end{itemize}
Hægt er að teikna rafrásir með TikZ, sjá mynd~\ref{fig:TikZCircuit}.
\begin{figure}[!ht]
  \caption{Einföld RL rafrás.}
  \begin{center}
    \begin{circuitikz}[american voltages, american currents, scale=0.50]
      \draw
        (0,0) to[short, -*] (5,0)
        to[short, *-* ] (10,0)
        to[short, *-* ] (15,0)
        to[short, *-  ] (20,0)
        ;
      \draw
        (0,0) to[I, l_=$\SI{20}{\ampere}$] (0,5)
        (5,0) to[R, l_=$\SI{0,5}{\ohm}$] (5,5)
        (10,0) to[L, l_=$\SI{5}{\henry}$, i^<=$i_L$] (10,5)
        (15,0) to[R, l_=$\SI{20}{\ohm}$, v^>=$\nu_0$, i_<=$i_0$, *-*] (15,5)
        (20,0) to[R, l_=$\SI{10}{\ohm}$ ] (20,5)
        ;
      \draw
        (0,5) to[short, -*] (5,5)
        (5,5) to[opening switch, -*, l_=$t{=}0$] (10,5)
        (10,5) to[short] (15,5)
        (15,5) to[short] (20,5)
        ;
    \end{circuitikz}
  \end{center}
  \label{fig:TikZCircuit}
\end{figure}
\subsubsection{Dæmi um texta úr rafrásargreiningu:}
	Látum $i$ vera fall einungis af tíma, þá getum við aðskilið jöfnurnar með
	\begin{align*}
		-R i &=  L \frac{di}{dt} && \Leftrightarrow \\
		-\frac{R}{L} dt &=  \frac{di}{i}
	\end{align*}
	með því að heilda báðar hliðar frá tímapunktinum $0^+$ til $t$ og straumi $i_0$ til $i$
	fæst
	\begin{align}
		\int_{i_0}^{i} \frac{dy}{y} &= -\frac{R}{L} \int_{0^+}^{t}  d\tau  
			&& \Leftrightarrow \nonumber \\
		\ln \left( \frac{i(t)}{i_0} \right) &= -\frac{R}{L} t
			&& \Leftrightarrow \nonumber \\
		i(t) &= i_0 e^{-\frac{R}{L} t} \label{equ:NaturalRLsolution}
\end{align}
þá er hægt að vísa í jöfnu~\ref{equ:NaturalRLsolution}.


%
% Fontes
%
\printbibliography[heading=bibintoc] % Adiciona fontes e traz uma tabela de conteudos
%escrever aqui


%
% Anexos
%
\appendix
\renewcommand{\chaptername}{Anexos}
\chapter{Anexos}
Dicas úteis e waypoints

\end{document}